\section{A strict inequality for norms}\label{sec:norm}

We now bound the higher spherical harmonics of a power of a norm by two
of its means.  Convexity supplies the differential estimate; the remaining
step consists of scalar inequalities.  Throughout this section set
\begin{equation}\label{eq:k-choice}
 k=\begin{cases}1,&n=4,\\2,&n\ge5,\end{cases}
 \qquad c=c_{n,k},\qquad d_* =\lambda_4-\lambda_2=2n+8.
\end{equation}
Thus the first power is used in dimension four and the square in higher
dimensions.  The scalar estimates below establish these two choices with
the same strict conclusion.

\begin{theorem}\label{thm:norm-estimate}
Let $g$ be a norm on $\R^n$, where $n\ge4$, normalized by
$\E g^{-n}=1$ on $\Sph$.  Then
\begin{equation}\label{eq:norm-estimate}
 c_{n,k}\|Q(g^k)\|_2^2
 \le (\E g^k)^2-(\E g^{k-1})^2.
\end{equation}
The inequality is strict unless $g=1$ on $\Sph$.
\end{theorem}

\subsection{Convexity and a scalar decomposition}

For $z>0$ define
\begin{align}
 e(z)&=\frac{z^{-n}-1}{n},&
 M(z)&=z^k-1+k e(z),&
 W(z)&=z^k-z^{k-1}+e(z).\label{eq:scalar-basic}
\end{align}
Since $M'(z)=k(z^{k-1}-z^{-n-1})$, the function $M$ decreases up to
$1$, increases thereafter, and satisfies $M(1)=M'(1)=0$.  In particular,
$M\ge0$.  For a norm with $\E g^{-n}=1$, the correction $e(g)$
has mean zero. We will split $g^k-1$ into a term controlled by
convexity and a remainder controlled by its variance. Choose the constants
\begin{equation}\label{eq:scalar-constants}
\begin{array}{c|cc}
 n&\beta&w\\ \hline
 4&1/2&3/14\\
 5,6&2/3&1/3\\
 n\ge7&2/3&1/2
\end{array}
\end{equation}
and split
\begin{equation}\label{eq:scalar-split}
 G(z)=\beta M(z)\mathbf 1_{\{z>1\}},\qquad
 D(z)=z^k-1-G(z).
\end{equation}
Thus $G$ vanishes for $z\le1$, while $D$ retains the remaining part
of $z^k-1$. To estimate $G(g)$ by convexity,
let $I(z)=0$ for $0<z\le1$, and, for $z>1$, let
\begin{align}
 I(z)&=\int_1^z M'(a)^2\,da\notag\\
 &=k^2\left(\frac{z^{2k-1}-1}{2k-1}
 -\frac{2(1-z^{k-1-n})}{n+1-k}
 +\frac{1-z^{-2n-1}}{2n+1}\right).
 \label{eq:scalar-I}
\end{align}
Both $G$ and $I$ are continuously differentiable across $1$;
$I\ge0$ and $G'^2=\beta^2 I'$.

\begin{lemma}\label{lem:curvature}
For every norm $g$ on $\R^n$,
\begin{equation}\label{eq:norm-curvature}
 2n\Var(G(g))+d_*\|QG(g)\|_2^2
 \le p\beta^2\E[gI(g)],\qquad p=n-1.
\end{equation}
\end{lemma}

\begin{proof}
Suppose first that $g$ is smooth away from the origin.  The trace of its
Euclidean Hessian on the tangent space at $u\in\Sph$ is
$\lapS g(u)+p g(u)$, and it is nonnegative by convexity.  Multiplying by
$I(g)\ge0$ and integrating by parts gives
\begin{align*}
 \E|\gradS G(g)|^2
 &=\beta^2\E[I'(g)|\gradS g|^2]
 =-\beta^2\E[I(g)\lapS g]
 \le p\beta^2\E[gI(g)].
\end{align*}
The spectral gap estimate \eqref{eq:spectral-gap} proves
\eqref{eq:norm-curvature} in this case.

Here is a passage to arbitrary norms that does not assume the
approximants are homogeneous.  Convolve the ambient norm with smooth,
even, nonnegative mollifiers of mass one and shrinking compact support,
and write the resulting smooth convex functions as $g_\varepsilon$.
On the unit sphere their tangent Hessian trace is
\[
 \lapS g_\varepsilon+p\partial_r g_\varepsilon\ge0.
\]
The preceding argument and \eqref{eq:spectral-gap} therefore give
\[
 2n\Var(G(g_\varepsilon))+d_*\|QG(g_\varepsilon)\|_2^2
 \le p\beta^2\E[I(g_\varepsilon)\partial_r g_\varepsilon].
\]
The functions $g_\varepsilon$ converge uniformly to $g$ on compact
sets, are uniformly Lipschitz there, and their gradients converge to
$\nabla g$ at every differentiability point of $g$.  To see the last
assertion, every limit of their bounded gradients satisfies the
subgradient inequality for $g$ by convexity and uniform convergence;
differentiability makes that subgradient unique.  The exceptional
directions have spherical measure zero, as in
Lemma~\ref{lem:gauge-variation}.  Homogeneity of $g$ then gives
$\partial_r g_\varepsilon(u)\to\nabla g(u)\cdot u=g(u)$ almost
everywhere.  On the sphere the approximants stay in a fixed positive
compact interval, while their radial derivatives are bounded.  Uniform
convergence on the left and dominated convergence on the right prove
\eqref{eq:norm-curvature}.
\end{proof}

\subsection{From scalar bounds to the norm estimate}\label{subsec:norm-deduction}

Define, for $z>0$,
\begin{equation}\label{eq:scalar-L}
 \mathcal L(z)=c\left(
 \frac{p\beta^2 zI(z)-2nG(z)^2}{wd_*}
 +\frac{D(z)^2}{1-w}\right).
\end{equation}

\begin{lemma}\label{lem:scalar-estimates}
The constants and functions above satisfy
\begin{equation}\label{eq:scalar-mean-bound}
 c\left(\frac{2n\beta^2}{wd_*}
       -\frac{(1-\beta)^2}{1-w}\right)
 \le\frac{2k-1}{k^2}
\end{equation}
and
\begin{equation}\label{eq:scalar-pointwise}
 \mathcal L(z)\le2W(z)\qquad(z>0).
\end{equation}
Equality in \eqref{eq:scalar-pointwise} holds only at $z=1$.
\end{lemma}


We first deduce the norm estimate from the two scalar bounds. Their
verification, including the strict pointwise inequality in every
dimension, occupies the remainder of the section.

\begin{proof}[Proof of Theorem~\ref{thm:norm-estimate}]
In the following expectations, $G,D,I$ denote their compositions with
$g$.  Set
\[
 m=\E g^k-1,\qquad \ell=\E g^{k-1}-1,\qquad t=\E G.
\]
The normalization gives $\E e(g)=0$, so $m=\E M(g)\ge0$,
$\E D=m-t$, and $\E W(g)=m-\ell$.  Since
$0\le G\le\beta M$, we have $0\le t\le\beta m$.
For every $a>0$, Jensen's inequality applied to the convex function
$y\mapsto y^{-a/n}$ yields
$\E g^a\ge(\E g^{-n})^{-a/n}=1$.
Together with
\[
 z^{k-1}\le\frac{(k-1)z^k+1}{k},
\]
this gives
\begin{equation}\label{eq:norm-mean-range}
 0\le\ell\le\frac{k-1}{k}\,m.
\end{equation}
For $k=1$ this simply says $\ell=0$.

Since $Q$ annihilates constants,
$Q(g^k)=QG+QD$.  The Hilbert-space inequality
\[
 \|a+b\|^2\le\frac{\|a\|^2}{w}
                       +\frac{\|b\|^2}{1-w}
 \qquad(0<w<1)
\]
and $\|QD\|_2^2\le\Var(D)$, followed by
Lemma~\ref{lem:curvature}, give
\begin{align}
 c\|Q(g^k)\|_2^2
 &\le\frac{c}{wd_*}
       \left(p\beta^2\E[gI(g)]-2n\Var(G)\right)
       +\frac{c}{1-w}\Var(D)\notag\\
 &=\E\mathcal L(g)+c\left(
       \frac{2n}{wd_*}t^2-\frac{(m-t)^2}{1-w}\right).
 \label{eq:norm-reduction}
\end{align}
The quadratic expression in parentheses is increasing for
$0\le t\le m$: its derivative is
\[
 \frac{4n}{wd_*}t+\frac{2(m-t)}{1-w}\ge0.
\]
It is therefore no larger than its value at $t=\beta m$.
Using \eqref{eq:scalar-mean-bound} and
\eqref{eq:norm-mean-range}, we conclude that the final term of
\eqref{eq:norm-reduction} is at most
\[
 \frac{2k-1}{k^2}m^2\le m^2-\ell^2.
\]
Now \eqref{eq:scalar-pointwise} yields
\begin{align*}
 c\|Q(g^k)\|_2^2
 &\le 2\E W(g)+m^2-\ell^2\\
 &=2(m-\ell)+m^2-\ell^2
   =(\E g^k)^2-(\E g^{k-1})^2,
\end{align*}
as claimed.  If $g$ is not identically $1$ on the sphere, continuity
gives a set of positive spherical measure on which $g\ne1$.
Lemma~\ref{lem:scalar-estimates} then gives
$\E\mathcal L(g)<2\E W(g)$, so the final inequality is strict.
The only constant norm on the sphere with $\E g^{-n}=1$ is $g=1$;
in this case both sides of \eqref{eq:norm-estimate} vanish.
\end{proof}

\subsection{Proof of the scalar bounds}

\begin{proof}[Proof of Lemma~\ref{lem:scalar-estimates}]
We first record the values of the spectral constant:
\begin{equation}\label{eq:scalar-c}
 c=\frac{121}{35}\quad(n=4),\qquad
 c=\frac{n^2+5n+7}{(n+1)(n+3)}\quad(n\ge5).
\end{equation}
For $n=4,5,6$ the left side of \eqref{eq:scalar-mean-bound} is,
respectively,
\[
 \frac{11}{12},\qquad \frac{589}{864},\qquad \frac{1387}{1890}.
\]
These are smaller than $1,3/4,3/4$, the respective right sides.
For $n\ge7$ the left side is
\[
 \frac{2(3n-4)(n^2+5n+7)}{9(n+1)(n+3)(n+4)}.
\]
Its difference from $3/4$, with the latter first, is
\[
 \frac{3n^3+128n^2+505n+548}{36(n+1)(n+3)(n+4)}>0.
\]
This proves \eqref{eq:scalar-mean-bound} for every dimension.

\medskip
\noindent\emph{The interval $0<z\le1$.}
Here $G=I=0$ and $|D(z)|=|z^k-1|\le k|z-1|$.
Also $W(1)=W'(1)=0$ and
\[
 W''(z)=\begin{cases}
 (n+1)z^{-n-2},&k=1,\\
 2+(n+1)z^{-n-2},&k=2.
 \end{cases}
\]
Thus $2W(z)\ge(n+2k-1)(z-1)^2$.  The positive margins
\begin{equation}\label{eq:scalar-local-margin}
 n+2k-1-\frac{ck^2}{1-w}
 =\begin{cases}
 3/5,&n=4,\\
 7/8,&n=5,\\
 43/21,&n=6,\\
 n+3-8c,&n\ge7
 \end{cases}
\end{equation}
give the desired strict inequality for $z<1$.  In the last case the
margin is $9/10$ at $n=7$ and increases thereafter: regarding $n$ as a
real variable in \eqref{eq:scalar-c},
\[
 c'(n)=-\frac{n^2+8n+13}{(n+1)^2(n+3)^2}<0.
\]
At $z=1$ both sides of \eqref{eq:scalar-pointwise} vanish.

\medskip
\noindent\emph{The interval $z>1$ in dimensions $4,5,6$.}
Set $x=z-1>0$.  Multiplication by a positive factor clears the negative
powers of $z$: if
\[
 P_n(z)=s_nz^{2n}\bigl(2W(z)-\mathcal L(z)\bigr),
 \qquad (s_4,s_5,s_6)=(2880,712800,221130),
\]
then substitution of \eqref{eq:scalar-basic}--\eqref{eq:scalar-L}
and collection of equal powers give
\begin{align}
 P_4(z)={}&462z^{10}-983z^9+93z^8
             +3399z^5-3378z^4+407,\notag\\
 P_5(z)={}&151525z^{14}-273570z^{12}-680800z^{11}
             +938157z^{10}\notag\\
         &\quad+225720z^7-462264z^5+101232,\notag\\
 P_6(z)={}&48399z^{16}-75894z^{14}-218004z^{13}
             +285961z^{12}\notag\\
         &\quad+56940z^8-127478z^6+30076.
 \label{eq:small-polynomials}
\end{align}
For clarity, the coefficients needed from $P_n(1+x)$ are listed below;
the constant and linear coefficients are zero.
\[
\begin{array}{c|rrr}
 [x^j]&n=4&n=5&n=6\\ \hline
 j=2&1728&623700&452790\\
 j=3&-1446&-1505900&679380\\
 j=4&-6711&-5802775&-132249\\
 j=5&1173&12849430&7800780\\
 j=6&17052&86314305&51269452\\
 j=7&20796&191506920&146058120
\end{array}
\]
All remaining coefficients are nonnegative, as can be seen without
listing them.  Let $a_nz^{d_n}$ be the leading term in
\eqref{eq:small-polynomials}.  At $j=8$, its contribution
$a_n\binom{d_n}{8}$ minus the contributions of all negative terms is
\[
 11943,\qquad 207280425,\qquad 114414300
 \quad\text{for } n=4,5,6,
\]
respectively.  For every $i<d_n$, the ratio
$\binom{i}{j}/\binom{d_n}{j}$ is nonincreasing in $j$; while it is
nonzero, its ratio at successive indices is $(i-j)/(d_n-j)\le1$.
Once $j>i$ it is zero.  Thus the same leading term covers all negative
contributions for every $8\le j\le d_n$.

It follows that $P_n(1+x)/x^2$ is bounded below by, respectively,
\begin{align*}
 U_4(x)&=1728-1446x-6711x^2+17052x^4,\\
 U_5(x)&=623700-1505900x-5802775x^2+86314305x^4,\\
 U_6(x)&=452790-132249x^2+51269452x^4.
\end{align*}
These lower bounds are positive.  Indeed
\begin{align*}
 1446x&\le400+1400x^2,&
 8111x^2&\le1000+17052x^4,\\
 1505900x&\le200000+2900000x^2,&
 8702775x^2&\le230000+86314305x^4.
\end{align*}
Each follows from $dt\le a+bt^2$ when $d^2\le4ab$; the corresponding
values of $4ab-d^2$ are
\[
 149084,\quad 2419679,\quad 52265190000,\quad 3670867899375.
\]
The first row gives $U_4(x)\ge328$ and the second gives
$U_5(x)\ge193700$.  Finally $U_6$, regarded as a quadratic in $x^2$,
has positive leading coefficient and negative discriminant, since
\[
 4\cdot452790\cdot51269452-132249^2=92839690886319>0.
\]
Therefore $P_n(z)>0$ for $z>1$ in all three dimensions.

\medskip
\noindent\emph{The interval $z>1$ in all dimensions $n\ge7$.}
Put $H=18(n+4)/c$.  Direct collection of powers in
\eqref{eq:scalar-L} now gives
\begin{equation}\label{eq:large-polynomial}
 \frac{9(n+4)}{c}z^{2n}(2W-\mathcal L)
 =z^{2n}(A_4z^4+A_2z^2+A_1z+A_0)
   +z^n(B_2z^2+B_0)+T_0,
\end{equation}
where
\begin{align}
 A_4&=\frac{2n-8}{3},&
 A_2&=H-12n-32-\frac{64}{n},\notag\\
 A_1&=-H+\frac{16(n-1)}3+32-\frac{16(n-1)}{2n+1},&&\notag\\
 A_0&=-\frac Hn+6n+40+\frac{64}{n}-\frac{128}{n^2},&&\notag\\
 B_2&=16+\frac{64}{n},&
 B_0&=\frac Hn-48-\frac{64}{n}+\frac{256}{n^2},\notag\\
 T_0&=\frac{16(n-1)}{2n+1}-\frac{128}{n^2}.&&
 \label{eq:large-coefficients}
\end{align}
We prove that the polynomial on the right of
\eqref{eq:large-polynomial}, expressed in $x=z-1$, has nonnegative
coefficients and a positive quadratic coefficient.

Write its two parentheses as $\sum_{i=0}^4q_ix^i$ and
$\sum_{i=0}^2r_ix^i$.  Thus
\begin{align*}
 q_0&=A_4+A_2+A_1+A_0,&q_1&=4A_4+2A_2+A_1,\\
 q_2&=6A_4+A_2,&q_3&=4A_4,\qquad q_4=A_4,\\
 r_0&=B_2+B_0,&r_1&=2B_2,\qquad r_2=B_2.
\end{align*}
In particular $q_3,q_4,r_1,r_2>0$.  The constant and linear
coefficients of the full polynomial vanish.  Its quadratic coefficient
is
\begin{equation}\label{eq:large-quadratic}
 \frac{9(n+4)}c\,(n+3-8c)>0.
\end{equation}
One can obtain these three coefficients either from
\eqref{eq:large-coefficients} or directly from the order of vanishing at
$z=1$: $G=O(x^2)$, $I=O(x^3)$, and $D=2x+O(x^2)$, so that
$2W-\mathcal L=(n+3-8c)x^2+O(x^3)$.  Positivity is
\eqref{eq:scalar-local-margin}.

The remaining signs admit explicit polynomial certificates.  Set
\begin{align*}
 V_0(n)&=34n^5-95n^4-1666n^3-5701n^2-8052n-2688,\\
 V_1(n)&=2n^5+9n^4-162n^3-845n^2-1220n-448,\\
 V_2(n)&=5n^3+8n^2-41n-56,\\
 R(n)&=7n^4+8n^3-187n^2-748n-896.
\end{align*}
Substituting $H=18(n+4)(n+1)(n+3)/(n^2+5n+7)$ gives
\begin{align}
 q_0+\frac{r_0}{8}
 &=\frac{V_0(n)}{4n^2(2n+1)(n^2+5n+7)},\notag\\
 q_1&=\frac{2V_1(n)}{n(2n+1)(n^2+5n+7)},\notag\\
 q_2&=\frac{2(n+4)V_2(n)}{n(n^2+5n+7)},&
 r_0&=-\frac{2R(n)}{n^2(n^2+5n+7)}.
 \label{eq:large-signs}
\end{align}
The required signs hold throughout their stated ranges because, for
$t\ge0$,
\begin{align}
 V_0(10+t)={}&34t^5+1605t^4+28534t^3
                  +227319t^2\notag\\
             &\quad+698128t+130692,\notag\\
 V_1(10+t)={}&2t^5+109t^4+2198t^3+19695t^2
                  +69280t+30852,\notag\\
 V_2(7+t)={}&5t^3+113t^2+806t+1764,\notag\\
 R(7+t)={}&7t^4+204t^3+2039t^2+7414t+4256.
 \label{eq:shift-certificates}
\end{align}
Consequently $q_2>0$ and $r_0<0$ for every $n\ge7$, while
$q_0+r_0/8>0$ and $q_1>0$ for every $n\ge10$.

Use the convention that $\binom mj=0$ when $j<0$ or $j>m$.  For
$3\le j\le2n$, comparison of the factors in the binomial products gives
\begin{equation}\label{eq:binomial-comparison}
 \binom nj\le 2^{-j}\binom{2n}{j}.
\end{equation}
Indeed $(n-a)/(2n-a)\le1/2$ for $0\le a<n$; if $j>n$ the left
side is zero.  For $n\ge10$, the sum of the $q_0$ and $r_0$
contributions to the coefficient of $x^j$ is therefore at least
\[
 \left(q_0+\frac{r_0}{2^j}\right)\binom{2n}{j}
 \ge\left(q_0+\frac{r_0}{8}\right)\binom{2n}{j}>0.
\]
All other contributions are nonnegative.  For $j>2n$, the $q_0$ and
$r_0$ contributions vanish and every remaining term is nonnegative.
Together with \eqref{eq:large-quadratic}, this proves coefficient
positivity for every $n\ge10$.

It remains to treat $n=7,8,9$, where $q_1$ can be negative.
For $3\le j\le2n$, divide the coefficient of $x^j$ by
$\binom{2n}{j}$, discard the nonnegative $q_3,q_4,r_2$ contributions,
and use \eqref{eq:binomial-comparison} for the negative $r_0$ term.
The resulting lower bound is
\begin{equation}\label{eq:small-large-dimension-bound}
 \begin{split}
 q_0+\frac{r_0}{2^j}
 +\frac{j}{2n-j+1}
   \left(q_1+q_2\frac{j-1}{2n-j+2}\right)
 +r_1\frac{\binom n{j-1}}{\binom{2n}{j}}.
 \end{split}
\end{equation}
Substitution in \eqref{eq:large-coefficients} gives the componentwise
lower bounds
\[
\begin{array}{c|rrrrr}
 n&q_0&q_1&q_2&r_0&r_1\\ \hline
 7&-2&-11&60&-2&50\\
 8&-1&-6&72&-6&46\\
 9&-1&-6&72&-6&46
\end{array}
\]
All weights in \eqref{eq:small-large-dimension-bound} are nonnegative,
so these bounds may be substituted directly.  They give the following
positive numbers at $j=3,4$; at $j=5$ we have already omitted the final
$r_1$ term:
\[
\begin{array}{c|ccc}
 n&j=3&j=4&j=5\ \text{without the final term}\\ \hline
 7&5/26&1233/1144&589/176\\
 8&37/28&10707/3640&1153/208\\
 9&101/136&3851/2040&395/112
\end{array}
\]
These finite substitutions can also be checked from the exact tuples
$(q_0,q_1,q_2,r_0,r_1)$:
\begin{align*}
 n=7:&\quad
 \left(-\frac{5984}{3185},-\frac{4832}{455},
       \frac{792}{13},-\frac{1216}{637},\frac{352}{7}\right),\\
 n=8:&\quad
 \left(-\frac{421}{629},-\frac{3736}{629},
       \frac{2688}{37},-\frac{145}{37},48\right),\\
 n=9:&\quad
 \left(\frac{2416}{10773},-\frac{2192}{1197},
       \frac{100568}{1197},-\frac{57968}{10773},\frac{416}{9}\right).
\end{align*}
For $5\le j\le2n$, the parenthesis
$q_1+q_2(j-1)/(2n-j+2)$ is positive and increasing in $j$.
In fact the displayed componentwise bounds give at $j=5$ the positive
lower bounds $119/11$, $210/13$, and $66/5$, respectively.
The prefactor $j/(2n-j+1)$ is positive and increasing too, while
$r_0/2^j$ increases because $r_0<0$.  The expression
\eqref{eq:small-large-dimension-bound} with its last term omitted is
therefore increasing throughout this range.  Positivity at $j=5$
proves positivity at every subsequent index through $2n$.

For $j>2n$ all $r_i$ contributions vanish, since $n+2<2n+1$.
The only possibly negative term is the $q_1$ contribution at
$j=2n+1$.  At this index the $q_1$ and $q_2$ contributions sum to
$q_1+2nq_2$, which the same componentwise bounds make at least
$829,1146,1290$ for $n=7,8,9$, respectively.  For $j\ge2n+2$ only
the nonnegative $q_2,q_3,q_4$ terms remain.  This handles every
coefficient, including the endpoints.

We have proved that the polynomial in \eqref{eq:large-polynomial}
has nonnegative coefficients and a strictly positive coefficient of
$x^2$ for every $n\ge7$.  It is consequently positive for every
$x>0$.  Its multiplier $9(n+4)z^{2n}/c$ is positive, so
$2W(z)-\mathcal L(z)>0$ for $z>1$.  This finishes the proof of
\eqref{eq:scalar-pointwise} and its strictness in every dimension.
\end{proof}

